\part*{第八部分：公开义务，不作软化}
\addcontentsline{toc}{part}{第八部分：公开义务，不作软化}
\stepcounter{partanchor}\setcounter{section}{0}

\section{确切地已确立的内容}

为免歧义，本卷宗所能支持的最强陈述是：

\begin{enumerate}[leftmargin=1.6em,itemsep=3pt]
\item \textbf{一个上界。} 存在一个显式的 $100$ 中心有理解，其覆盖半径（covering radius）至多为 $U_\star=1154469548855829/10^{16}=0.1154469548855829$，由
$138\,437$ 个精确整数对象认证（$128\,436$ 个内部三元组、$9\,900$ 个边界对、$100$
个对径对、$1$ 个例外情形），并带有精确的覆盖平方半径界
$1332799939235381342323027277241/10^{33}=U_\star^2$，把参考值
$115446954885583/10^{15}$ 恰好改进 $1/10^{16}$。
\item \textbf{一个下界。} $B31$--$B35$ 在 $R=2831/25000$ 处由精确证书（certificate）排除；
$B\ge37$ 在 $R=2831/25000$ 处以截距（intercept）$19/2500$ 排除；$B=36$ 在
$R_\star=22647/200000$ 处以截距 $939/125000$ 在 $1\,841\,089$ 个有理盒上排除。半径
单调性随后同样给出 $R_\star$ 处的 $B31$--$B35$ 与 $B\ge37$，故
\[
\rD(100)>22647/200000=0.113235 .
\]
\item \textbf{不依赖构型的界。} $\rD(100)\ge2787/25000=0.11148$（长弦），并有
$\ge1111720153/10^{10}$（加强包络）与 $\ge554570693/5000000000$（整数边计数）
同样成立。
\item \textbf{一个精确的障碍。} 在冻结的 capped-pair 常数下，$B=36$ 松弛在
$R=2831/25000$ 处于有理点 $x=y=9921033/12500000$ 上同时超过旧截距与面积所允许的最大截距。这是关于松弛的陈述，而不是关于覆盖的陈述。
\end{enumerate}

其余一切都或者是条件性的、诊断（diagnostic）性的、接口层面的，或者是公开的。

\section{公开义务}

{\footnotesize
\begin{longtable}{@{}>{\raggedright\arraybackslash}p{0.30\textwidth}>{\raggedright\arraybackslash}p{0.58\textwidth}@{}}
\toprule
义务（obligation） & 缺失内容的精确陈述 \\
\midrule
\endfirsthead
\toprule
义务（obligation） & 缺失内容的精确陈述 \\
\midrule
\endhead
\bottomrule
\endlastfoot
$\rD(100)$ 的精确值 & 未获得。被接受的区间宽度约为 $0.0022$；没有任何论证把它闭合。\texttt{global\_\allowbreak{}optimality\_\allowbreak{}proved = false} \\
$\rD(100)$ 的极小多项式 & 未获得：不是系数，不是次数，不是不可约性，也不是有理的根隔离区间。存在的是\emph{某一个构型}覆盖半径的极小多项式
$67t^2-1$，那是一个不同的对象，并且被明确地如此划定范围 \\
闭合成单个代数数 & 两个端点是不相同的数；没有任何东西把理想包络的零点与最优半径等同起来，而包络的零点明确地是关于该方法极限的陈述，而不是关于 $\rD(100)$ 的陈述 \\
完整的内部映射枚举 & 唯一已知的必要计数定义域是被接受的归一化几何归约下得到的 $87$ 个 $(m,B)$ 对，其中
$m=96..100$（$96:B32..35$；$97:B29..42$；$98:B28..46$；$99:B28..50$；$100:B28..54$）。它不是完整的拓扑枚举，而且存活下来的族中仍含有不可行的图 \\
$m<100$ 与其他图 & 计数桥在其被重放的各半径处排除了 $m\le99$，但对 $m<100$、对其他图族以及对退化的普遍处理尚未闭合 \\
退化 & 重合的中心与冗余的圆盘由归一化处理；其余的退化目录（相切接触、共线中心、边界多重活跃）在各个分支（branch）内部由显式论证处理，而不是由单一的、统一的机器检查处理 \\
$B=36$ 不等式 & 直接的配对/角度松弛无法闭合 $0.11324$；一个共同端点的透镜约束或一个真正耦合的多边平均不等式已被认定为缺失的成分，但它没有通过的全局证书 \\
可续跑的尾部基础设施 & $0.113236$ 与 $0.113237$ 的尝试是停在墙钟时间上，而不是停在数学上；一个带缓存的逆三角与支撑区间、并按对角/非对角分片的可续跑盒队列已被设计出来，但尚未交付 \\
精确值后端 & 被吸收的平方松弛 / 实量词消去理论是唯一已知的通向精确值的可终止路线，但它没有 $n=100$ 的后端；被编译系统的零点集处处奇异，因此常规的 Krawczyk 盒无法替代完整的 QE 后端，而变量预算（$605$ 个几何 $+$ $1318$ 个松弛，$1922$ 个被消去）是一个预算，而不是一个运行时间 \\
遗留源码激活 & 全部 $174$ 个遗留条件源码返回
\texttt{PRECONDITION\_\allowbreak{}MISSING}；依赖门是 \texttt{UNDECIDED}。「每一个历史变体都可直接调用并被逐一重新计时」\emph{未}达成 \\
输出隔离欠账 & 计时、组合/缩放与共享支撑这三个入口点仍然缺少统一的、可重复输出的包装器；整合报告把这列为尚欠的工作 \\
几何形式化 & 不存在 Lean/Coq 内核验证。每一条全称陈述都是解析散文加精确有限算术；解析/几何之间的桥是一个经人工审阅的依赖，在每份证书中都记录为 \texttt{machine\_\allowbreak{}verified\_\allowbreak{}geometry = false} \\
未发布的常数 & 至少有一个尾部常数是从一个不属于镜像的包导入的；它的值在本地被重新表述并数值地重现（附录~\ref{app:budget}），但把该常数与 $3\sqrt3/4$ 的一个向外包络等同起来是一个推断 \\
记号欠账 & 报告没有给出计数参数 $M$ 的自足定义；
它只通过整数超出恒等式与存活剖面出现。读者必须从证书中把它重建出来，而本专著指出这一点，而不是发明一个定义 \\
\end{longtable}}


\section{被撤回或被更正的断言}

一份研究记录只有在保留其撤回时才可信。以下断言曾被提出，随后被撤回或被更正，而这些更正也是记录的一部分：

\begin{itemize}[leftmargin=1.4em,itemsep=3pt]
\item \textbf{$4\arcsin(r/2)$ 弧界。} 错误；$\partial D\cap B(c_i,r)$ 的正确最大角测度是
$2\arcsin r$，在 $\|c_i\|=\sqrt{1-r^2}$ 处取到。错误的常数
在 $r=1$ 附近无法推出 $<\!\pi$，因而不可用。
\item \textbf{边界环带「全局无增益」的接受。} 已作废：一个错误的内面过滤器使待定集为空，而独立脚本复制了同一个错误。退出码 $0$
不是覆盖域正确的证据。
\item \textbf{stage-47 v2 对全部 $4776$ 个构型的排除。} 已否决。
\item \textbf{角能量惯性断言 $+25/-6$。} 已撤回：对真实数据做 CAS 得到 $32\times32$ 角坐标中的一维负方向（外加一个零 = 全局旋转），投影到 $31$ 维间隙坐标后仍是一维负方向。不再作任何惯性断言。
\item \textbf{「径向壳 LP 给出接近 $92.56$ 的已认证覆盖数界。」}
已更正：采样解在 $4001$ 点细网格上不可行（覆盖 $0.99903<1$，
对偶 $1.00317>1$）；唯一被认证的陈述是 $r=L$ 处的粗括界
$[1/r^2,\ \pi/\mathrm{lens}(1,r)]=[80.911,\,165.734]$。更强的说法「这条路线已关闭／证书不可能」也被撤回。
\item \textbf{「构型一侧仍有 $5$--$6\%$ 的余地。」} 已撤回：小 $n$ 的最优解是环型族，而最好的 $n=100$ 构型是晶格型族，因此它们不可比较；晶格族的比值 $c(n)$ 为 $2.13$（$n=40$）、$1.92$（$50$）、$1.35$（$70$）、
$1.14$（$100$）、$1.12$（$140$）、$0.82$（$200$）。
\item \textbf{stage-142 的发现层数字}（耳近似 $1.332634$ 对真值
$1.42016417919847$，以及一处 $K$-目标勘误）不得作为证书值引用。
\item \textbf{stage-44 的第一版}（$4496$ 次收缩 / $597$ 次关联）已被
$2248/562$ 取代；\textbf{stage-45/46} 使用了错误的外 dart 所有者；依赖审计（audit）中一处历史脚本哈希不匹配是被\emph{诊断}出来的，而不是被悄悄修复的。
\end{itemize}

\section{对差距的诚实解读}

两个数字勾勒出局面。不依赖构型的最好下界是 $0.11148$，它比引擎所接受的界 $0.113235$ 低
$1.55\%$；而被接受的区间是
$[0.113235,\,0.1154469548855829]$，相对宽度约为 $1.92\%$。因此引擎比任何单个可用的解析不等式都强得多，而剩下的工作不是更多精度或更多盒子的
问题：它是针对一个分支、在一个半径上的一个不等式，外加 --- 对于精确值而言 --- 一个在 $n=100$ 处尚不存在的枚举与消去后端。

同样重要的是说明当前的界\emph{不}对什么敏感。增加盒子、提高截距、复用 $B\ge37$ 的截距，或者在冻结族内同步重调斜率与罚系数，都会因可证明的理由而失败（第五部分）。该界也不受算力限制：$B=36$ 的候选运行只花 $130$ s 的报告时间，而整个引擎中最大的单项开销是尾部细分。限制它的是：当前局部配对表达式在障碍点（obstruction point）处确实大于 $B=36$ 面积预算所能支付的量。
